\documentclass{article}
\usepackage{graphicx} % Required for inserting images
\usepackage{cancel}
\usepackage{ulem}
\usepackage{amsmath} % Required for math environments
\usepackage{xcolor}
\title{An Overview of String Theory}
\author{Arpita Dhar(2105137)}

\begin{document}

\maketitle

\section{Introduction}
\st{\underline{\textbf{string}}} “String theory” proposes that point-like particles of particle physics are
replaced by one-dimensional objects called strings [1]. It describes how these
strings propagate through space and interact with each other.

\section{The Basics of String Theory}

In string theory, the basic idea is that the fundamental constituents of reality
are strings of incredibly small scale (approximately $10^{-35}$ meters in length), 
and their vibrations determine the mass and force charges of the particles. See
Figure 1 for visual representations of string theory concepts.

\section{Theoretical Dimensions}
Table 1 presents different theoretical dimensions and their implications in string
theory.
\begin{center}
\caption{Table 1: Various String Theory Models}
\begin{tabular}{ |p{10em}| p{7em}| p{14em}|}
  \hline
  Models & Dimensions & Implications \\
  \hline
  Row 1, Col 1 & Row 1, Col 2 & Row 1, Col 3 \\
  \hline
  Row 2, Col 1 & Row 2, Col 2 & Row 2, Col 3 \\
  \hline
  Row 3, Col 1 & Row 3, Col 2 & Row 3, Col 3 \\
  \hline
\end{tabular}
\end{center}

\footnote{Intentionally done for \LaTeX\ skill evaluation.}

\pagebreak

\section{Dynamics of Strings}
The Nambu-Goto action in string theory, describing the dynamics of a classical
string, is given by: \\
\[
S = \int  \sqrt{-\det(g_{ab})} \, d\tau \, d\sigma
\]
where: \\
• \( S \) is the action integral.\\
• \( T \) is the string tension.\\
• \( g_{ab} \) is the induced metric on the worldsheet of the string.\\
• \( \tau \) (tau) and \( \sigma \) (sigma) are the parameters on the worldsheet.\\

\vspace{2em}
\underline{Note:} \textcolor{red}{ This article is solely designed for evaluating students’ \LATEX\ skills.}.

\section{References}

\bibliographystyle{ieetr}
\bibliography{ref}
\end{document}
